% Source: https://artofproblemsolving.com/community/c3h2936505_2017_emcc_guts_round_p13_p24__exeter_math_club_competition
% Stable identity: aops-topic-2936505-q18
Both rectangles have diagonal $BD$, hence share its midpoint $O$ and circumcircle. Their other diagonals are $AC$ and $EF$. Since $ACEF$ forms a square, these two diameters are perpendicular. Congruence means the acute angles which $BD$ makes with $AC$ and $EF$ are equal, so both are $45^\circ$. If the common diagonal length is $d$, either rectangle has area $\dfrac{1}{2}d^2\sin45^\circ=\dfrac{d^2}{2\sqrt2}=1$. The square has diagonals $d,d$, giving area $\dfrac{d^2}{2}=\boxed{\sqrt2}$.
\begin{asy}
usepackage("amsmath");
size(7cm);
pair A=dir(0),C=dir(180),Ee=dir(90),F=dir(270),B=dir(45),D=dir(225),O=(0,0);draw(A--B--C--D--cycle,blue);draw(B--Ee--D--F--cycle,red);draw(A--Ee--C--F--cycle,dashed);draw(A--C^^Ee--F^^B--D,dashed);label("$A$",A,E);label("$B$",B,NE);label("$C$",C,W);label("$D$",D,SW);label("$E$",Ee,N);label("$F$",F,S);label("$O$",O,S);
\end{asy}
